\justifying
\NSFCBodyText

\textbf{（1）可学习乘积量化与状态空间模型协同的长上下文高效注意力机制研究。}资助来源为 CAAI—蚂蚁科研基金，项目编号202604140006，获资助总额15万元，执行期2026年2月至2027年2月。申请人为参与者，正式负责高效注意力机制研究。本申请针对运行时任务状态、证据恢复和预算配置，主要固定执行模型；两项工作分别研究模型内部计算效率与智能体外部上下文管理，不将前一项目任务或成果重复列为本项目新交付。

\textbf{（2）面向政企数字化应用的智能知识库与问答系统研发。}委托单位为新疆数字云创有限公司，横向项目编号202604140013，合同经费12万元，执行期2026年6月至2026年10月，申请人主持并按合同构建一个知识库。该项目提供知识组织及问答系统研发背景；本申请拟研究可复位多步骤任务中的压缩、恢复和全开销关系。既有项目成果及知识产权由合同双方共同享有，本申请不预设其代码、数据可公开或直接复用。

\textbf{（3）多模态融合技术研究及其在新疆产业集群中的应用。}资助来源为自治区天池英才引进计划（青年博士人才），无独立项目批准号，以任务书为依据；经费50万元，执行期2025年2月至2028年2月，申请人主持并负责总体研究方案、关键模型设计与成果组织。该项目提供多源表示和实验组织基础，本申请的核心任务为文本与工具上下文管理，不将多模态应用任务重复安排为本项目交付。

\textbf{（4）融合多源数据大模型驱动的光伏发电功率预测方法研究。}资助来源为自治区高校基本科研业务费科研项目培育类，批准号XJEDU2026P005，经费5万元（直接经费3.9万元、间接经费1.1万元），执行期2026年1月至2026年12月，申请人主持，负责总体方案、模型研究与成果组织。该项目以功率预测为对象，与本申请的多步骤任务证据管理在研究对象与评价终点上不同，仅提供多源模型与实验管理经验。

以上内容依据已确认项目事实列示。主要参与者及其在研项目情况随\textbf{【待补 F-KS-02：本项目人员及相关在研项目清单】}一并核对，正式申报时依据企业基金表格确定最终列示范围，不据此推断全部参与者均无其它项目。
